\subsection{LoRA replication}
\label{sec:lora}
LoRA learns a weaker shortcut at rank 4 ($G =$ \gapLoraRFour) and a stronger one at rank 64 (\gapLoraRSixtyFour), with
clean twins near zero (\gapLoraRFourClean; \gapLoraRSixtyFourClean). On the find side the picture is the same as for
full fine-tuning: on \emph{opposite} reviews the lens reads the true sentiment at layer~9 (rank 4: \lensCalOppLNineLoraRFour;
rank 64: \lensCalOppLNineLoraRSixtyFour), and swapping in clean blocks 9--11 removes nearly the whole gap (rank 4:
\swapARemovedLateLoraRFour; rank 64: \swapARemovedLateLoraRSixtyFour). However, control (c-ii) \emph{fails} under LoRA
(fraction kept: rank 4 \swapCiiKeptLateLoraRFour; rank 64 \swapCiiKeptLateLoraRSixtyFour; required $\ge 0.80$), so the
LoRA swaps cannot be read as clean removals.

On the fix side, rank 64 behaves like full fine-tuning: every single block removes the shortcut, again because of the
counter-examples (Sec.~\ref{sec:counterexample-ablation}). Rank 4 is the exception: only \ccPassWithNamesLoraRFour{} of
its 12 single blocks remove it (block~1, for example, removes \retrainRemovedBOneLoraRFour). Retraining always updates
whole blocks, so this reflects the rank-4 shortcut itself, not a limited capacity of the fix. Swap--retrain agreement
holds for rank 4 (Spearman \spearmanSwapRetrainLoraRFour) but not for rank 64 (\spearmanSwapRetrainLoraRSixtyFour).

\textbf{Why does control (c-ii) fail under LoRA?} A LoRA model differs from another seed's model in its adapter
initialization and in its training path (data order and dropout). We trained rank-64 models on a $3\times3$ grid of
initialization seeds and data seeds (one of the nine collapsed during training and was excluded by a rule fixed before
any swap) and swapped blocks 9--11 between pairs that differ in only the data seed (A), only the initialization (B), or
both (C). Over all pairs, swaps fail to keep the shortcut in \lFiveFailPairsCellA{} pairs of cell A and
\lFiveFailPairsCellC{} of cell C, but in \lFiveFailPairsCellB{} of cell B. Descriptively, the training path rather than
the initialization makes rank-64 blocks incompatible across runs, in line with the view that independently trained
networks need not share a representation basis. The registered test used three replicate pairs per cell (fraction kept
A \lFiveKeptCellA, B \lFiveKeptCellB, C \lFiveKeptCellC); with roughly one failing pair in three it read the original
failure as a single-pair event (Table~\ref{tab:prereg}).

\textbf{Where LoRA is trained (Phase L4).} If adapters are placed only on blocks 0--5, the name's effect still reaches
the answer position late, at layer \lFourArrivalEarly{} (first layer with $\Fl \ge 0.5$), so the late read-out is largely
a property of GPT-2's own later blocks, which here are frozen. With adapters only on blocks 6--11 it arrives at
\lFourArrivalLate, so the arrival also moves with where the shortcut is learned. The early-only shortcut is weaker
(gap \lFourGapEarly, against \lFourGapLate{} for the late-only set).

\textbf{A partly reliable name (Phase L6).} To separate rank from shortcut strength, we tried to give rank 64 a
shortcut as weak as rank 4's by making the name agree with the label in only 90\%, 80\%, 70\% or 60\% of the named
reviews. The seed-0 gaps were \lSixGapPNinetysZero, \lSixGapPEightysZero, \lSixGapPSeventysZero{} and
\lSixGapPSixtysZero, so no level matched rank 4 and the comparison was not run. Descriptively (one seed), the curve
itself is informative: at 90\% reliability, only 50 contradicting reviews out of 10{,}000 reduce the gap from
\gapLoraRSixtyFoursZero{} to \lSixGapPNinetysZero. A few counter-examples present during training already weaken the shortcut
sharply, mirroring their role in the repair data (Sec.~\ref{sec:counterexample-ablation}).
